\documentclass[a4paper, 12pt, oneside, headspline]{scrartcl}
% !TeX root = ../expose.tex
% Variablen welche innerhalb der gesamten Arbeit zur Verfügung stehen sollen
\newcommand{\titleDocument}{Exposé}
\newcommand{\subjectDocument}{Vergleich von Python Test-Tools für Test-driven development}
\newcommand{\institution}{Hochschule für angewandte Wissenschaften Augsburg}
\newcommand{\insSection}{Fakultät für Informatik}
\newcommand{\name}{Maximilian Konter}
\newcommand{\email}{maximilian.konter@hs-augsburg.de}



% Farbige links
\usepackage{color}

% Umlaute unter UTF8 nutzen
\usepackage[utf8]{inputenc}

% deutsche Silbentrennung
\usepackage[ngerman]{babel}

% Grafiken aus PNG Dateien einbinden
\usepackage{graphicx}

% bricht lange URLs "schoend" um
\usepackage[hyphens,obeyspaces,spaces]{url}

% mehrseitige Tabellen ermöglichen
\usepackage{longtable}

% für Tabellen
\usepackage{array}

% Schaltet den zusätzlichen Zwischenraum ab, den LaTeX normalerweise nach einem Satzzeichen einfügt.
\frenchspacing

% Paket für Zeilenabstand
\usepackage{setspace}

% für Bildbezeichner
\usepackage{capt-of}

% für Stichwortverzeichnis
\usepackage{makeidx}

% Festlegung Art der Zitierung - Havardmethode: Abkuerzung Autor + Jahr
\bibliographystyle{alphadin}

% für Listings
\usepackage{listings}
\lstset{numbers=left, numberstyle=\tiny, numbersep=5pt, keywordstyle=\color{black}\bfseries, stringstyle=\ttfamily,showstringspaces=false,basicstyle=\footnotesize,captionpos=b}
\lstset{language=python}

% Indexerstellung
\makeindex

% Abkürzungsverzeichnis
\usepackage[german]{nomencl}
\let\abbrev\nomenclature

% Abkürzungsverzeichnis LiveTex Version
\renewcommand{\nomname}{Abkürzungsverzeichnis}
\setlength{\nomlabelwidth}{.25\hsize}
\renewcommand{\nomlabel}[1]{#1 \dotfill}
\setlength{\nomitemsep}{-\parsep}
\makenomenclature

% Disable single lines at the start of a paragraph (Schusterjungen)
\clubpenalty = 10000
% Disable single lines at the end of a paragraph (Hurenkinder)
\widowpenalty = 10000
\displaywidowpenalty = 10000

% Hyperlinks im Text
\usepackage[bookmarksnumbered,pdftitle={\titleDocument},hyperfootnotes=false]{hyperref}
\hypersetup{colorlinks, citecolor=red, linkcolor=blue, urlcolor=black}

% neue Kopfzeilen mit fancypaket
\usepackage{fancyhdr} %Paket laden
\pagestyle{fancy} %eigener Seitenstil
\fancyhf{} %alle Kopf- und Fußzeilenfelder bereinigen
\fancyhead[L]{\nouppercase{\leftmark}} %Kopfzeile links
\fancyhead[C]{} %zentrierte Kopfzeile
\fancyhead[R]{\thepage} %Kopfzeile rechts
\renewcommand{\headrulewidth}{0.4pt} %obere Trennlinie
\fancyfoot[C]{\thepage} %Seitennummer
\renewcommand{\footrulewidth}{0.4pt} %untere Trennlinie

\begin{document}	
	% Titelseite
	% !TeX root = ../bachelor.tex
\thispagestyle{empty}

\begin{figure}[t]
 \centering
 \includegraphics[width=0.3\textwidth]{abb/logoHSA}
\end{figure}


\begin{verbatim}


\end{verbatim}

\begin{center}
\Large{\institution}
\end{center}


\begin{center}
\Large{\insSection}
\end{center}
\begin{verbatim}




\end{verbatim}
\begin{center}
\doublespacing
\textbf{\LARGE{\titleDocument}}
\singlespacing
\begin{verbatim}

\end{verbatim}
\textbf{{~\subjectDocument~}}
\end{center}
\begin{verbatim}

\end{verbatim}
\begin{center}

\end{center}
\begin{verbatim}

\end{verbatim}
\begin{center}
\textbf{zur Erlangung des akademischen Grades \\ Bachelor of Science}
\end{center}
\begin{verbatim}





\end{verbatim}

\begin{footnotesize}
\begin{flushleft}
\begin{tabular}{llll}
\textbf{Thema:} & & \subjectDocument & \\
& & \\
\textbf{Autor:} & & \name & \\
& & \email & \\
& & MatNr. 951004 & \\
& & \\
\textbf{Version vom:} & & \today &\\
& & \\
\textbf{1. Betreuer:} & & Dipl.-Inf. (FH), Dipl.-De Erich Seifert, MA &\\
\textbf{2. BetreuerIn:} & & Prof. Dr.-Ing. Alexandra Teynor &\\
\end{tabular}
\end{flushleft}
\end{footnotesize}
	
	% Inhaltsverzeichnis anzeigen
	\newpage
	\tableofcontents
	
	% das Abkürzungsverzeichnis
	\newpage
	% Abkürzungsverzeichnis soll im Inhaltsverzeichnis auftauchen
	\addcontentsline{toc}{section}{Abkürzungsverzeichnis}
	% das Abkürzungsverzeichnis entgültige Ausgeben
	\fancyhead[L]{Abkürzungsverzeichnis} %Kopfzeile links
	% !TeX root = ../../main.tex
\nomenclature{GPL}{GNU General Public License}
\nomenclature{GNU}{GNU is not Unix}
\nomenclature{LGPL}{GNU Lesser General Public License}
\nomenclature{TDD}{Test-driven development}
\nomenclature{GUI}{Graphical User Interface}
\nomenclature{STDLIB}{Python Standard Bibliothek}
\nomenclature{VCS}{\gls{vcs}}
	\printnomenclature
	
	%%%%% Beginn Dokument %%%%%
	\newpage
	\fancyhead[L]{\nouppercase{\leftmark}} %Kopfzeile links
	
	% 1,5 facher Zeilenabstand
	\onehalfspacing
	
	% einzelne Kapitel
	\section{Ausgangslage}
Was ist der aktuelle stand bezüglich TDD?

\section{Problembeschreibung}
\begin{itemize}
	\item Warum werden tests immer wichtiger und beliebter?
	\item Wieso nutzen Firmen nicht TDD sondern lieber andere Methoden?
	\item Testen dauert - $Das Problem soll durch die automation geloest werden$
	\item weitere
\end{itemize}

\section{Fragestellung}
Wieso sollte man TDD machen?
	
	% !TeX root = ../expose.tex
\section{Zielsetzung}
Hier beschreibe ich was ich mit dieser Arbeit erreichen will und wieso.
	% !TeX root = ../expose.tex
\section{Theoretische Grundlage}
Welche literatur gibt es?
	% !TeX root = ../expose.tex
\section{Konzept}
Hypothese: TDD ist besser?
Methodik: Wie soll dies bewiesen werden? Anhand der Planung einer komplette umgebung
	% !TeX root = ../expose.tex
\section{Vorläufige Gliederung}
\begin{enumerate}
	\item Abstract
	\item Einleitung
\end{enumerate}
	% !TeX root = expose.tex
\section{Motivation}
Warum will ich darueber scheiben?
	
	\onecolumn
	
	% einfacher Zeilenabstand
	\singlespacing
	
	% Literaturliste soll im Inhaltsverzeichnis auftauchen
	\newpage
	\addcontentsline{toc}{section}{Literaturverzeichnis}
	
	% Literaturverzeichnis anzeigen
	\renewcommand\refname{Literaturverzeichnis}
	\bibliography{expose}
\end{document}